\section{Evaluación del cumplimiento de requisitos}
Una vez finalizada la ejecución de los casos de prueba, se realizó la evaluación del cumplimiento de los requisitos establecidos para la solución. Para ello, se contrastaron los resultados obtenidos y las evidencias recopiladas durante las pruebas con las condiciones definidas para cada requisito.

El resultado de esta evaluación se clasificó en tres categorías: cumple, cuando las funcionalidades requeridas fueron implementadas y verificadas mediante las pruebas; cumplimiento parcial, cuando se implementó una parte del requisito, pero alguna de sus características no pudo ser incorporada dentro de la arquitectura establecida; y no cumple, cuando la funcionalidad requerida no fue implementada en la solución. Los resultados consolidados se presentan en la \autoref{tab:evaluacion-requisitos}.

\begingroup
\fontsize{8.5pt}{9.5pt}\selectfont
\renewcommand{\arraystretch}{1.3}
\setlength{\tabcolsep}{4pt}

\begin{longtable}{|
	>{\raggedright\arraybackslash}p{1.2cm}|
	>{\raggedright\arraybackslash}p{4.0cm}|
	>{\raggedright\arraybackslash}p{2.3cm}|
	>{\raggedright\arraybackslash}p{2.0cm}|
	>{\raggedright\arraybackslash}p{4.0cm}|}

	\caption{Evaluación del cumplimiento de requisitos}
	\label{tab:evaluacion-requisitos}
	\\
	\hline
	\tableheader
	\textbf{ID}                &
	\textbf{Requisito}         &
	\textbf{Pruebas asociadas} &
	\textbf{Resultado}         &
	\textbf{Observaciones}
	\\
	\hline
	\endfirsthead

	\caption{Evaluación del cumplimiento de requisitos (continuación)}
	\\
	\hline
	\tableheader
	\textbf{ID}                &
	\textbf{Requisito}         &
	\textbf{Pruebas asociadas} &
	\textbf{Resultado}         &
	\textbf{Observaciones}
	\\
	\hline
	\endhead

	R-01                       & Gestionar centralizadamente las identidades y mantener autorización diferenciada por sistema.                                          & PT-01, PT-02                  & Cumple               & Se verificó la autenticación centralizada y la aplicación de autorización diferenciada mediante grupos.                                                                         \\ \hline
	R-02                       & Gestionar el ciclo de vida de las cuentas, incluyendo creación, actualización, habilitación, deshabilitación, revocación y expiración. & PT-03                         & Cumple               & Se verificaron las operaciones de gestión del ciclo de vida contempladas en la implementación.                                                                                  \\ \hline
	R-03                       & Definir y aplicar políticas de seguridad para usuarios y credenciales.                                                                 & PT-04                         & Cumple               & Se verificó la aplicación de las políticas de credenciales configuradas.                                                                                                        \\ \hline
	R-04                       & Controlar y monitorear los accesos y facilitar la revisión de permisos.                                                                & PT-02, PT-06                  & Cumple               & Se verificó el control de acceso mediante grupos y la revisión de los permisos asociados.                                                                                       \\ \hline
	R-05                       & Registrar actividades relevantes de gestión de identidades y accesos para fines de auditoría.                                          & PT-05                         & Cumple               & Se verificó que una operación de gestión de identidad generara un registro que permite identificar la actividad realizada.                                                      \\ \hline
	R-06                       & Fortalecer la autenticación mediante políticas de credenciales y mecanismos adicionales de autenticación.                              & PT-04                         & Cumplimiento parcial & Se implementaron políticas de credenciales, pero la integración de \acs{MFA} con los servicios mediante \acs{LDAP} no fue posible dentro de la arquitectura implementada.       \\ \hline
	R-07                       & Proporcionar mecanismos seguros de recuperación o restablecimiento de credenciales.                                                    &                               & No cumple            & FreeIPA no proporciona un mecanismo nativo de recuperación de contraseña mediante autoservicio. La incorporación de alternativas externas quedó fuera del alcance del proyecto. \\ \hline
	R-08                       & Permitir la evolución y ampliación del servicio ante la incorporación de nuevos sistemas o activos.                                    & PT-01, pruebas de integración & Cumple               & La integración de los tres servicios mediante el directorio centralizado evidencia la capacidad de incorporar sistemas al servicio de directorio.                               \\ \hline
\end{longtable}
\endgroup
\normalsize


\subsection{Consideraciones y limitaciones}
Durante la ejecución de las pruebas se tuvieron en cuenta las condiciones del entorno de implementación y el alcance establecido para la solución. Estas consideraciones permitieron delimitar los resultados obtenidos y establecer las condiciones bajo las cuales se evaluó el cumplimiento de los requisitos.

Las pruebas se realizaron sobre el entorno compuesto por el servidor principal, el servidor réplica de FreeIPA y el balanceador HAProxy, junto con los servicios integrados OpenVPN, OpenProject y Xen Orchestra.

Para la ejecución de las pruebas se utilizaron tanto la interfaz web de administración de FreeIPA como la interfaz de línea de comandos. La interfaz web se empleó principalmente para evidenciar elementos cuya representación gráfica facilita su interpretación, como la topología de replicación. Por su parte, la \ac{CLI} se utilizó para las operaciones y consultas que requerían contrastar directamente la información entre el servidor principal y el servidor réplica, especialmente en las pruebas de replicación, continuidad y recuperación.

Las evidencias se integraron directamente con el procedimiento de cada caso de prueba. De esta manera, cada evidencia se relaciona con la actividad que permitió obtenerla, facilitando la trazabilidad entre las acciones ejecutadas y los resultados observados. Asimismo, cuando una actividad ya había sido validada en una prueba anterior, sus resultados se utilizaron como antecedente en pruebas posteriores para evitar la repetición innecesaria de procedimientos y evidencias.

\subsubsection{Limitaciones}
Una de las limitaciones identificadas corresponde al requisito R-06, relacionado con el fortalecimiento de la autenticación mediante políticas de credenciales y mecanismos adicionales de autenticación. La autenticación multifactor se contempló dentro del alcance del proyecto; sin embargo, la integración de OpenVPN, OpenProject y Xen Orchestra se realizó mediante \ac{LDAP}, y las implementaciones \ac{LDAP} utilizadas por estos servicios no contemplan, dentro de este mecanismo de integración, el uso de la autenticación multifactor proporcionada por FreeIPA. Por tanto, aunque FreeIPA dispone de capacidades relacionadas con mecanismos adicionales de autenticación, estas no pudieron incorporarse a los tres servicios mediante la arquitectura implementada.

La incorporación de autenticación multifactor directamente en los servicios podría requerir mecanismos de integración adicionales o componentes específicos de cada plataforma, lo que implicaría ampliar la arquitectura planteada y utilizar servicios adicionales. Por este motivo, dicha alternativa se mantiene fuera del alcance de la implementación realizada. En consecuencia, la evaluación de R-06 debe distinguir entre las políticas de credenciales implementadas y el mecanismo adicional de autenticación que no pudo integrarse mediante \ac{LDAP}.

Otra limitación corresponde al requisito R-07, relacionado con la recuperación o restablecimiento de credenciales. La implementación permite al usuario modificar su contraseña cuando dispone de la credencial vigente; sin embargo, no proporciona un mecanismo de recuperación de contraseña cuando esta ha sido olvidada. Esta limitación corresponde a una decisión del propio proyecto FreeIPA: su documentación señala que el restablecimiento de contraseña mediante autoservicio no forma parte de las funcionalidades principales de gestión de usuarios y que los usuarios que olvidan sus credenciales deben recurrir al administrador o a la mesa de ayuda, previa verificación de su identidad \citep{FreeIPAReset2026}.

La documentación de FreeIPA también señala que determinados mecanismos habituales de recuperación, como las preguntas de seguridad o el restablecimiento mediante correo electrónico, pueden introducir vulnerabilidades asociadas, por ejemplo, a ingeniería social o a la transmisión insegura de información. Debido a estas consideraciones, el equipo de FreeIPA no contempla incorporar este tipo de mecanismos dentro de las funcionalidades principales de la plataforma \citep{FreeIPAReset2026}; en su lugar, plantea que los administradores que requieran esta capacidad desarrollen e integren soluciones propias de restablecimiento mediante la \ac{API} interna o la interfaz \ac{LDAP} del sistema, adaptadas a los requerimientos de seguridad de cada infraestructura.
